\chapter{Fundamentação Teórica}

\section{Aprendizado Contrastivo e Modelos Fundacionais em Ciência dos Materiais}
O aprendizado contrastivo baseia-se na maximização da concordância mútua entre diferentes representações ou modalidades de uma mesma amostra no espaço latente. Dado um conjunto de dados, o objetivo é aprender uma função de mapeamento que projete amostras semanticamente semelhantes para regiões próximas no espaço latente, enquanto amostras distintas são projetadas para regiões distantes. A formalização inicial deste paradigma foi introduzida por Hadsell et al.~\cite{hadsell2006dimensionality}, que propuseram a função de perda contrastiva (\textit{contrastive loss}) baseada em pares de atração e repulsão por margem. Com o avanço do aprendizado profundo, essa abordagem evoluiu para o aprendizado contrastivo auto-supervisionado com o arcabouço SimCLR, proposto por Chen et al.~\cite{chen2020simple}, que utiliza transformações estocásticas (aumentações) e a perda NT-Xent (\textit{Normalized Temperature-scaled Cross Entropy}) para maximizar a concordância entre diferentes visões do mesmo dado sem a necessidade de rótulos. Complementarmente, Khosla et al.~\cite{khosla2020supervised} estenderam esse conceito para o escopo supervisionado (\textit{Supervised Contrastive Learning} - SupCon), permitindo que múltiplos exemplos que compartilham o mesmo rótulo ou faixa de propriedade atuem conjuntamente como pares positivos no espaço de representação.

Na ciência dos materiais, essa abordagem tem se consolidado no pré-treinamento de codificadores (\textit{encoders}) integrados a modelos fundacionais (\textit{foundation models}), os quais são treinados com grandes volumes de dados não rotulados para posteriormente serem ajustados (\textit{fine-tuning}) em tarefas específicas de regressão ou classificação de propriedades~\cite{van2026survey}.
Essa metodologia evoluiu para arquiteturas multimodais capazes de integrar representações cristalográficas, composicionais e estruturais. Um exemplo proeminente é o modelo MultiMat, que estende arquiteturas do tipo CLIP (\textit{Contrastive Language-Image Pre-training}) para o domínio dos materiais~\cite{moro2025multimodal}. Nesse paradigma, codificadores neurais dedicados processam cada modalidade individual do material, projetando as informações em um espaço latente unificado. Combinado a arquiteturas avançadas, como Transformers e redes convolucionais 3D, esse aprendizado representacional mitiga a dependência de grandes conjuntos de dados rotulados, viabilizando a predição precisa de propriedades complexas, como bandgap, tensor elástico e módulo de compressão (\textit{bulk modulus})~\cite{moro2025multimodal}.

\section{Representação por Grafos e Transformações Estruturais em Sistemas Cristalinos}

A aplicação eficiente do aprendizado contrastivo a materiais cristalinos e bidimensionais exige a conversão das estruturas atômicas em representações matemáticas que respeitem a invariância rotacional, translacional e a simetria de permutação atômica. Nesse âmbito, as Redes Neurais em Grafos (GNNs), exemplificadas pela \textit{Crystal Graph Convolutional Neural Network} (CGCNN) desenvolvida por Xie e Grossman~\cite{xie2018crystal}, tornaram-se o padrão para modelar materiais periódicos, convertendo átomos em nós e ligações interatômicas em arestas ponderadas.

No âmbito do aprendizado contrastivo auto-supervisionado, a etapa de aumento de dados (\textit{data augmentation}) desempenha um papel crítico ao introduzir perturbações geométricas e topológicas na estrutura do grafo sem alterar suas propriedades fundamentais. Arquiteturas como o CrystalCLR~\cite{koker2022graph} demonstraram que a aplicação de transformações físicas diretas na estrutura cristalina permite que o modelo aprenda representações latentes invariantes significativamente mais expressivas do que descritores (\textit{fingerprints}) projetados manualmente. O pré-treinamento desses modelos sobre vastas bases de dados, como o C2DB~\cite{haastrup2018computational} para sistemas bidimensionais e o Materials Project para cristais tridimensionais, cujas propriedades de referência são calculadas via DFT por pacotes como o VASP~\cite{kresse1996efficient}, provê um arcabouço computacional escalável para superar o gargalo do elevado custo computacional~\cite{langer2022representations} na descoberta de materiais 2D.
